\documentclass[10pt,a4paper]{amsart}
\usepackage[margin=19mm]{geometry}
\usepackage{amsmath,amssymb,mathtools}
\usepackage[scr=rsfs,cal=euler]{mathalfa}
\usepackage{xfrac}
\usepackage[unicode,hidelinks,bookmarks=false]{hyperref}
\newcommand{\ie}{i.e.\ }
\newcommand{\cf}{cf.\ }
\newcommand{\resp}{respectively\ }
\newcommand{\Subsection}{\subsection}
\providecommand{\nobreakdash}{\nobreak\hbox{-}}
\setlength{\emergencystretch}{2em}
\multlinegap0pt
\usepackage{bm}
\usepackage{bbm}
\usepackage{dsfont}
\usepackage{xspace}
\usepackage{cancel}
\usepackage{mathtools}
\usepackage{amsthm}
\makeatletter
\@ifundefined{theorem}{\newtheorem{theorem}{Theorem}}{}
\@ifundefined{lemma}{\newtheorem{lemma}[theorem]{Lemma}}{}
\makeatother
\newcommand{\C}{\mathbb C}
\newcommand{\N}{\mathbb N}
\newcommand{\Q}{\mathbb Q}
\newcommand{\R}{\mathbb R}
\newcommand{\Z}{\mathbb Z}
\title[Deep zero problems and the HRT conjecture]{Deep zero problems and the HRT conjecture}
\author{Yufei Li, Zeguang Liu,, Kehe Zhu}
\date{Source: arXiv:2601.09080v1 (CC BY 4.0)}
\begin{document}
\maketitle

\noindent\emph{Source and licence.} Deep zero problems and the HRT conjecture. Version arXiv:2601.09080v1 (CC BY 4.0), distributed under the Creative Commons Attribution 4.0 International licence, as recorded in the arXiv licence metadata for this exact version and shown on the abstract page https://arxiv.org/abs/2601.09080v1. Full rights notice: licenses/arxiv-2601.09080-CC-BY-4.0.txt.\par\smallskip

\noindent\textit{Editorial scope.} Counted unit: the Theorem whose source label is \texttt{t4} in the article (source0.tex lines 419--428, source label \texttt{t4}), delivered together with the article's own complete original proof of it (source0.tex lines 430--472, one \texttt{proof} environment of 43 lines). No mathematics is written, repaired or re-derived: statement and proof are the article's own text line for line. Required context retained uncounted: t1 (lines 304--308); l2 (lines 366--368); l3 (lines 377--379). Every definition, display and label of those lines is retained unchanged. Omitted from this package and not redistributed: 1--303, 309--365, 369--376, 380--418, 429--429, 473--730. Nothing inside a retained excerpt is omitted or altered: every retained body is delivered exactly as the article's lines are written, so no omission receipt of any kind accompanies this package. Citations: the retained excerpts carry no citation command at all, so no bibliography entry of the article is transcribed and no reference is redistributed. Reference adaptation: none was needed and none was made; every reference inside the delivered unit resolves inside it, so no unresolved reference or citation is produced. Printed slips kept literal: no printed word, symbol, hypothesis or number of the counted statement or of its proof was altered. Layout and class: the article's own class is \texttt{amsart} with packages including stix2, amsmath, amsthm, amscd, amssymb, latexsym, upref, stmaryrd, enumitem, graphicx, tikz; the wrapper is the lane's pinned \texttt{amsart} editorial wrapper at 10pt on A4, which restates the theorem-like environments the retained text needs and adds the article's own macro definitions that the retained text needs (\texttt{\textbackslash C}, \texttt{\textbackslash N}, \texttt{\textbackslash Q}, \texttt{\textbackslash R}, \texttt{\textbackslash Z}), with extra wrapper packages (bm, bbm, dsfont, xspace, cancel, mathtools). The delivered numbering is the unit's own. Assets: no retained span contains a figure environment, a graphics-inclusion command, a PDF-page inclusion, a table or a TikZ picture; no image file is copied into this package, the package declares no asset dependency and no image file is redistributed with it. Licence: version arXiv:2601.09080v1 of this article is distributed under the Creative Commons Attribution 4.0 International licence, as recorded in the arXiv licence metadata for this exact version and shown on the abstract page https://arxiv.org/abs/2601.09080v1. A full rights notice is delivered at licenses/arxiv-2601.09080-CC-BY-4.0.txt. Characters: every retained excerpt is pure ASCII, so the delivered engine, XeLaTeX with its default Latin Modern text font, reports no missing character.
\begin{theorem}\label{t1}
Let $d\in\{1,2,3,4,6\}$ and $\beta\in\C\setminus\{0\}$. Then (the Fock 
space version of) the HRT conjecture is valid for the points
$$\lambda_k=\beta e^{2k\pi i/d},\qquad 0\le k\le d-1.$$
\end{theorem}

\begin{lemma}\label{l2}
$\phi(n)$ is even for all $n>2$. 
\end{lemma}

\begin{lemma}\label{l3}
The solutions to the equation $\phi(n)=2$ are $n=3, 4, 6$.
\end{lemma}

\begin{theorem}\label{t4}
Suppose $d\in\N$ and $\omega=e^{2\pi i/d}$. Then the
following are equivalent.
\begin{enumerate}
\item[(a)] The points
$$\lambda_k=\omega^k=e^{2k\pi i/d},\qquad 0\le k\le d-1,$$
belong to a regular lattice on the complex plane.
\item[(b)] $n\in\{1,2,3,4,6\}$.
\end{enumerate}
\end{theorem}

\begin{proof}
It is obvious that the points in $\{\omega^k: 0\le k\le d-1\}$ 
belong to a regular lattice in $\C$ when $d\in\{1,2,3,4\}$. 
By Theorem~\ref{t1}, this is also true when $d=6$.

Next we fix $d>4$ and assume that 
$\{\lambda_0,\lambda_1,\cdots,\lambda_{d-1}\}$ belong 
to a regular lattice
$$L=\Z e_1+\Z e_2+e$$ 
in the complex plane, where $e\in\C$ and $e_1,e_2$ are two 
complex numbers that are linearly independent with respect 
to $\R$.

For $1\le k\le d-1$ we write
$$v_k=\lambda_k-\lambda_0=n_{k1}e_1+n_{k2}e_2,$$
where $n_{k1}$ and $n_{k2}$ are integers. Then for $k\ge2$
we can write
$$\begin{pmatrix}v_1\\ v_k\end{pmatrix}=
\begin{pmatrix}n_{11} & n_{12}\\ n_{k1} & n_{k2}\end{pmatrix}
\begin{pmatrix}e_1\\ e_2\end{pmatrix}.$$
With respect to the real numbers, $v_1$ and $v_k$ are linearly
independent, and $e_1$ and $e_2$ are linearly independent. It
follows that each of the above $2\times2$ matrices of integers
is invertible. In particular, we have
$$\begin{pmatrix}e_1\\ e_2\end{pmatrix}=\begin{pmatrix}
n_{11} & n_{12}\\ n_{21} & n_{22}\end{pmatrix}^{-1}
\begin{pmatrix}v_1\\ v_2\end{pmatrix}.$$
Thus for $k\ge3$ we can write
$$\begin{pmatrix}v_1\\ v_k\end{pmatrix}=\begin{pmatrix}
n_{11} & n_{12}\\ n_{k1} & n_{k2}\end{pmatrix}
\begin{pmatrix}n_{11} & n_{12}\\ n_{21} & n_{22}\end{pmatrix}^{-1}
\begin{pmatrix}v_1 \\ v_2\end{pmatrix}.$$
Elementary linear algebra tells us that the inverse of an invertible
matrix with integer entries has rational entries. It follows that each
$v_k$, $k\ge3$, is a linear combination of $v_1$ and $v_2$ with
rational coefficients. Consequently, the sets 
$$\{v_0, v_1,\cdots, v_{d-1}\}\quad{\rm and}\quad
\{\lambda_0,\lambda_1,\cdots,\lambda_{d-1}\}$$ 
are both contained in ${\rm Span}_{\Q}\{1,\omega,\omega^2\}$,
which implies $\left[\Q(\omega) : \Q\right]\le3$. This along with
Lemma~\ref{l2} and Theorem E shows that $\phi(d)=2$. An
application of Lemma~\ref{l3} then gives $d=6$.
\end{proof}

\end{document}
